\documentclass[11pt]{article}
\usepackage[margin=1in]{geometry}
\usepackage{amsmath,amssymb,amsthm}
\usepackage[hidelinks]{hyperref}
\newtheorem{theorem}{Theorem}
\newtheorem{proposition}[theorem]{Proposition}
\newcommand{\nb}{\mathrm{nb}}
\newcommand{\sr}{\mathrm{s}}
\newcommand{\kh}{K_{\mathrm{hit}}}
\newcommand{\kp}{K_{\mathrm{proj}}}
\newcommand{\ks}{K_{\mathrm{SRW}}}
\newcommand{\one}{\mathbf 1}
\newcommand{\ip}[2]{\langle #1,#2\rangle}
\newcommand{\var}{\operatorname{Var}}
\newcommand{\ran}{\operatorname{ran}}
\newcommand{\tr}{\operatorname{tr}}
\newcommand{\E}{\mathbb E}
\setlength{\emergencystretch}{1em}
\title{Vertex hitting times for non-backtracking walks\\
       with fixed and stationary starts}
% Supply author names, affiliations and corresponding-author contact at submission.
\author{}
\date{}
\makeatletter
\renewcommand{\maketitle}{%
  \begin{center}
    {\LARGE\@title\par}
    \ifx\@author\@empty\else
      \vspace{0.6em}{\large\begin{tabular}[t]{c}\@author\end{tabular}\par}
    \fi
  \end{center}\vspace{0.4em}}
\makeatother
\begin{document}
\maketitle
\begin{abstract}
Suppressing immediate reversals can increase the mean time to hit a prescribed
neighbour. A cycle on 23 vertices with one chord gives mean hitting times
$13$ for non-backtracking walk and $1843/143$ for simple random walk,
contradicting a pointwise comparison of Altisen, Devismes, Gerbaud and
Lafourcade. On theta graphs the ratio can grow without bound. For a starting
vertex sampled with probability proportional to its degree, we prove the opposite inequality for
every fixed target and characterize equality, with forced reversal at leaves.
We also compare the two vertex Kemeny constants of Breen, Kempton, Knudson
and Shumway. Their difference is nonnegative and vanishes precisely when
the conditional mean return time at each vertex is independent of the first
step. The supplementary Lean~4 development verifies these results and the
two-sided bounds for the difference of the vertex Kemeny constants.
\end{abstract}
\noindent\textit{Keywords.} Non-backtracking walk, hitting time, Kemeny's constant,
tabu walk, graph Laplacian. \quad \textit{MSC 2020.} 05C81, 60J10.

\section{Two ways to choose a starting vertex}
Let $G=(V,E)$ be a finite connected simple undirected graph with at least two
vertices, $m=|E|$, $d_i=\deg(i)$ and $\pi_i=d_i/(2m)$. From a prescribed
vertex, both walks choose the first neighbour uniformly. Simple random walk
(SRW) uses this rule at every step. The non-backtracking walk (NB) chooses
uniformly among neighbours other than the preceding vertex, with forced
reversal at a leaf. Write $H^\bullet_{ia}=\E_i T_a$ for the mean first
hitting time of $a$, where $T_a=\inf\{n\ge0:X_n=a\}$ and
$\bullet\in\{\nb,\sr\}$. Define
\[
 t^\bullet_a=\sum_i\pi_i H^\bullet_{ia},\qquad
 \kh=\sum_a\pi_a t^\nb_a,\qquad \ks=\sum_a\pi_a t^\sr_a .
\]
Here $\ks$ is the usual Kemeny constant of SRW. An initial visit counts as
a hit, so the term with $i=a$ is zero in both averages.

Altisen et al.~\cite[Cor.~5, author version]{ADGL} assert that their
one-entry FIFO tabu walk reaches any specified target at least as fast in
expectation as any rule with memory of length zero or one. The zero-entry
rule is SRW. We exhibit an ordered pair for which this inequality is reversed.
The reversal persists on a family of maximum degree three, with an
unbounded ratio of the two expectations.

\begin{theorem}\label{thm:start}
For every $a\in V$, one has $t^\nb_a\le t^\sr_a$. Equality holds precisely
when every neighbour of $a$ is a leaf. Consequently $\kh<\ks$ unless
$G=K_2$, in which case the walks coincide.
\end{theorem}
The initial law in this theorem is essential. Both the reversal at a fixed
start and the strict stationary-start improvement occur for the \emph{same}
target in the example below. Fasino, Tonetto and Tudisco developed directed-edge
equations for second-order hitting times and compared uniform-start averages
on three empirical networks \cite[Sec.~6]{FTT}. Neal's comparison for
walks that suppress reversals concerns the asymptotic variance of
time averages along a stationary trajectory \cite{Neal}. Breen et al.~\cite{BKKS}
ask for the ordering of the degree-weighted vertex constants.
Theorem~\ref{thm:start} answers Question~4 of their preprint.
Section~\ref{sec:trace} gives a converse to their equality criterion
(Question~2) and bounds their gap (Question~3). Edge-space Kemeny constants
use an oriented edge as the target \cite{BFG}.

\section{A counterexample with adjacent start and target}\label{sec:theta}
Join $A$ and $B$ by three internally disjoint paths of lengths $1,p,q$,
where $p\ge3$ and $q\ge2$. Denote this graph by $\Theta(1,p,q)$. On the
$p$-path the first two vertices after $A$ are $a$ and $x$, in this order.
The walker starts at $x$ and aims for $a$. Figure~\ref{fig:witness} shows
the case $(p,q)=(8,15)$.

\begin{figure}[htbp]
\centering
\setlength{\unitlength}{1pt}
\begin{picture}(246,63)(-19,-2)
\qbezier(4,24)(104,66)(204,24)
\qbezier(4,24)(104,-18)(204,24)
\put(4,24){\line(1,0){200}}
\put(4,24){\circle*{3.8}}
\put(204,24){\circle*{3.8}}
\put(29,33.19){\circle*{2.8}}
\put(54,39.75){\circle*{2.8}}
\put(79,43.69){\circle*{2.8}}
\put(104,45){\circle*{2.8}}
\put(129,43.69){\circle*{2.8}}
\put(154,39.75){\circle*{2.8}}
\put(179,33.19){\circle*{2.8}}
\put(17.33,18.77){\circle*{2.8}}
\put(30.67,14.29){\circle*{2.8}}
\put(44,10.56){\circle*{2.8}}
\put(57.33,7.57){\circle*{2.8}}
\put(70.67,5.33){\circle*{2.8}}
\put(84,3.84){\circle*{2.8}}
\put(97.33,3.09){\circle*{2.8}}
\put(110.67,3.09){\circle*{2.8}}
\put(124,3.84){\circle*{2.8}}
\put(137.33,5.33){\circle*{2.8}}
\put(150.67,7.57){\circle*{2.8}}
\put(164,10.56){\circle*{2.8}}
\put(177.33,14.29){\circle*{2.8}}
\put(190.67,18.77){\circle*{2.8}}
\put(29,33.19){\circle{5.8}}
\put(54,39.75){\circle{5.8}}
\put(0,25){\makebox(0,0)[r]{\scriptsize $0$}}
\put(208,25){\makebox(0,0)[l]{\scriptsize $8$}}
\put(26,39){\makebox(0,0)[b]{\scriptsize $a=1$}}
\put(54,46){\makebox(0,0)[b]{\scriptsize $x=2$}}
\end{picture}
\caption{The 23-vertex witness. The upper and lower paths have eight and
fifteen edges. The chord joins $0$ and $8$.}\label{fig:witness}
\end{figure}

\begin{proposition}\label{prop:theta}
For this ordered pair,
\[
 H^\nb_{xa}=\frac{4p+3q+1}{6},\qquad
 H^\sr_{xa}=p-2+\frac{(2q+1)p+3q^2+4q+2}{(p+1)q+p}.
\]
For $(p,q)=(8,15)$, $H^\nb_{xa}-H^\sr_{xa}=16/143>0$.
Along $q=p^2$, $H^\nb_{xa}/H^\sr_{xa}\sim p/8$.
\end{proposition}
The graph is a cycle with one chord. For the witness and the asymptotic
family, the NB expression also follows from \cite[App.~A.1, case~(2)]{BKKS}:
take $n=p+q$, $k=p+1$ and label the $p$-path from $B=1$ to $A=k$.
The start and target then have labels $p-1$ and $p$, respectively.
We give both calculations to make the pointwise comparison self-contained.

\begin{proof}
Write $s=(1+q)/2$ for the mean length of a uniformly chosen bypass between
$A$ and $B$ (the single edge and the $q$-path are the bypasses).
For NB, the first step hits $a$ with probability $1/2$. In the other
direction it takes $p-2$ steps from $x$ to $B$ and then a mandatory,
uniformly chosen bypass to $A$. At each subsequent visit to a terminal
the walker either takes the target path or uses the other bypass, with
equal probabilities. The number $N$ of additional bypass crossings has
$\Pr(N=j)=2^{-j-1}$, so $\E N=1$ and $\Pr(N\text{ even})=2/3$.
Conditional on $N$, each bypass crossing has mean length $s$: the
initial bypass is uniform and subsequent bypasses alternate. The terminal
leg to $a$ has length $1$ for even $N$ and $p-1$ for odd $N$. Thus
\[
 2H^\nb_{xa}=1+(p-2)+2s+(p+1)/3 .
\]

For SRW, ground $a$ and inject a unit current at $x$, giving voltages
$V_i$ in the unit-resistance network. If $L_a$ is the grounded
combinatorial Laplacian, then $L_a V=e_x$ and
$L_a H^\sr_{\cdot a}=(d_i)_{i\ne a}$ by the first-step equations.
Symmetry of $L_a$ yields $H^\sr_{xa}=\sum_i d_i V_i$.
The bypasses have parallel resistance $t=q/(q+1)$. The detour from
$x$ to $a$ through $B$ and $A$ has resistance $p-1+t$, in parallel
with the direct edge $xa$. Its current is $I=1/(p+t)$, hence
\[
 V_A=I,\qquad V_B=(1+t)I,\qquad V_x=(p-1+t)I .
\]
Each edge contributes the sum of its endpoint voltages to $\sum_i d_iV_i$.
On a corridor of $\ell$ edges the voltages interpolate linearly, so the
corridor contributes $\ell$ times the sum of its two terminal voltages.
Adding the corridors gives
\[
 H^\sr_{xa}=(p-1)V_x+(p+q-1)V_B+(q+2)V_A.
\]
Substitution gives the stated formula. When $q=p^2$, the two means are
asymptotic to $p^2/2$ and $4p$, respectively.
\end{proof}

For the numerical witness take the cycle $0,1,\ldots,22,0$ and add
$\{0,8\}$. Thus $A=0$, $B=8$, $x=2$ and $a=1$. The graph has
24 edges, two vertices of degree three and 21 of degree two.
Section~2 of \cite{ADGL} starts the walk with an empty list. Its update
rule in Section~3 makes $\mathrm{FIFO}_1$ store exactly the preceding
vertex. The fallback when all neighbours are taboo never occurs here.
Every graph is 1-free in the terminology of \cite[Sec.~4]{ADGL}.
Hence this case satisfies the graph, memory and initialization hypotheses
of Corollary~5. The positive-time formulation in \cite[Cor.~9]{TR}
gives the same hitting time because $x\ne a$.

The two first directions for NB have conditional mean lengths $1$ and
$25$. In the electrical calculation $(V_A,V_B,V_x)=(16,31,127)/143$,
giving $(7\cdot127+22\cdot31+17\cdot16)/143=1843/143$.
For comparison, the exact first-step systems on this graph give the
following two rows, both with target $a=1$:
\[
\begin{array}{lcc}
 \text{initial vertex}&\text{NB}&\text{SRW}\\ \hline
 2 &13&1843/143\\[2pt]
 \text{drawn from }\pi&1087/72&51907/858
\end{array}
\]
The path argument in \cite[Sec.~6]{TR} removes retractions from a walk
stopped at its target and identifies the resulting path law with the FIFO
law. This identification already fails on a triangle $x,z,a$. Before SRW
hits $a$ it alternates between $x$ and $z$. After erasing immediate
reversals the path is $x,a$ with probability
$\sum_{k\ge0}2^{-(2k+1)}=2/3$. The fresh NB walk takes $x,a$ with
probability $1/2$. The triangle tests the path-law step. The strict
hitting-time counterexample is Proposition~\ref{prop:theta}.

\section{A positive identity for stationary starts}\label{sec:energy}
Index the $2m$ oriented-edge states by $(i,j)$, with $i$ the current vertex
and $j$ its predecessor. On real functions of these states use the uniform
probability inner product. Let $P$ be orthogonal projection onto functions
of $i$ alone (averaging over $j\sim i$), let $Q=I-P$ and set
$(Jw)(i,j)=w(j,i)$. We identify $\ran P$ with vertex functions in
$L^2(\pi)$. The operator $J$ is a self-adjoint isometry.

On $\ran Q$ let $R$ multiply by $d_i-1$. Its inverse is well defined:
the fibre of $\ran Q$ at a leaf is zero. Write $S=P-R^{-1}Q$.
On each fibre of degree at least two, $S$ has zero diagonal and off-diagonal
entries $1/(d_i-1)$. At a leaf it is the identity. The transition on
functions of the edge state is $F=SJ$. A fictitious predecessor chosen
uniformly conditional on $i$ reproduces the uniform first step, including
at a leaf.

The full edge chain is doubly stochastic. Each closed communicating class
projects onto every vertex: if it contains $(i,j)$, a predecessor in the
class has current vertex $j$, and every other neighbour of $i$ can occur
as the next current vertex. Its projected vertex set is therefore closed
under adjacency. Connectedness and finite-state recurrence imply that the
mean time to hit any vertex is finite from every edge state.

Use the block operators
\[
 L=I-PJP\ \text{on }\ran P,\qquad B=QJP:\ran P\to\ran Q,\qquad
 C=R+QJQ\ \text{on }\ran Q .
\]
Both $L$ and $C$ are positive semidefinite. Indeed
$\ip z{Lz}=(2m)^{-1}\sum_{\{i,j\}\in E}(z_i-z_j)^2$ for vertex
functions, and $R\ge I$ on $\ran Q$. The kernel of $L$ consists of constants.

Fix $a$. Let $w(i,j)$ be the NB mean time to hit $a$ from the edge state,
with $w=0$ over $a$. Put $u=Pw$, $v=Qw$ and
$b=\one-\one_{\{i=a\}}/\pi_a$. Thus $u_i=H^\nb_{ia}$. The first-step
equation $w=\one+Fw$ away from $a$ gives the full-space identities
\begin{equation}\label{eq:block}
 Lu-B^*v=b,\qquad Bu+Cv=e,\qquad e=\one_{\{i=a\}} QJw .
\end{equation}
The left-hand side of the first equation equals $1$ off $a$ and has
global mean zero because $L\one=B\one=0$. For the second identity, apply
$Q$ to the recurrence and multiply by $R$ off $a$. On the target fibre
$v=0$. If $r_{aj}$ denotes the mean positive return time to $a$ conditional
on first stepping to $j$, then $r_{aj}=1+w(j,a)$. Consequently
\begin{equation}\label{eq:variance}
 \|e\|^2=\pi_a\,\var_{j\ {\rm uniform\ in}\ N(a)}(r_{aj}).
\end{equation}

\begin{proof}[Proof of Theorem~\ref{thm:start}]
For $h_i=H^\sr_{ia}$ the ordinary recurrence gives $Lh=b$ and $h_a=0$.
Since $u_a=v(a,\cdot)=0$ and $e$ is supported over $a$,
\eqref{eq:block} implies $t^\nb_a=\ip b u=\ip u{Lu}+\ip v{Cv}$.
Now expand the first term on the right in
\begin{equation}\label{eq:first}
 t^\sr_a-t^\nb_a=\ip{h-u}{L(h-u)}+\ip v{Cv}\ \ge0 .
\end{equation}
If equality holds, then $h=u$ since both vanish at $a$, and $Cv=0$.
Equation~\eqref{eq:block} would imply $Bh=0$ off $a$. If a neighbour
$i$ of $a$ has another neighbour $k$, then $h_a=0<h_k$, so centring
the neighbour values at $i$ gives $(Bh)(i,\cdot)\ne0$. Conversely, if
every neighbour of $a$ is a leaf, connectedness makes $G$ a star with
centre $a$. Both walks hit $a$ in one step from every other vertex.
When $|V|\ge3$ some vertex has a neighbour of degree at least two.
Averaging over $a$ gives the strict comparison of the constants.
\end{proof}

\section{The gap between two vertex Kemeny constants}\label{sec:trace}
In this section assume $\delta=\min_i d_i\ge2$ and that $G$ is not a
cycle. The full edge chain is then irreducible \cite[Lemma~3.1]{OW}.
Let $W$ project onto the global constants and $Z=(I-F+W)^{-1}$.
This inverse also exists for an irreducible periodic chain. Let
$Z_v=PZP|_{\ran P}$ and put $\kp=\tr(Z_v)-1$. In the forward-edge
convention of \cite[Def.~3.6]{BKKS}, the edge transition has matrix
$JS=F^\top$ in the port indices used here. Each compressed diagonal
entry is the sum over the corresponding $d_a$-by-$d_a$ block divided by
$d_a$, which is invariant under transpose. Thus this $\kp$ agrees with
their projected fundamental-matrix constant.

\begin{theorem}\label{thm:trace}
One has $\kp\le\kh$, with equality exactly when $r_{aj}$ is independent
of $j$ at each $a$. Put $\Delta=\max_i d_i$ and
$V_r=\sum_a\pi_a^2\,\var_{j\ {\rm uniform\ in}\ N(a)}(r_{aj})$.
Let $\ell$ be the greatest length of a maximal chain whose internal
vertices have degree two and whose endpoints have degree at least three,
allowing both endpoints to coincide. Then
\[
 \frac{V_r}{\Delta}\ \le\ \kh-\kp\ \le\ 2\ell V_r .
\]
When $\delta>2$, the upper bound $V_r/(\delta-2)$ also holds.
\end{theorem}

\begin{proof}
For the same target $a$ as in \eqref{eq:block}, take $w_0=Zb$ plus a
constant chosen so that $u_0(a)=0$, with $u_0=Pw_0$, $v_0=Qw_0$.
The full Poisson equation $(I-F)w_0=b$ gives
$Lu_0-B^*v_0=b$ and $Bu_0+Cv_0=0$.
Since $Z\one=\one$ and $Z$ preserves global means, expansion of
$\ip b{Zb}$ gives $\ip b{u_0}=(Z_v)_{aa}/\pi_a-1$.
For $\xi=u-u_0$ and $\eta=v-v_0$, subtraction yields
$L\xi=B^*\eta$ and $B\xi+C\eta=e$. Self-adjointness gives
\[
 \ip b\xi=\ip{u_0}{L\xi}-\ip{v_0}{B\xi}
         =-\ip{v_0}{C\eta+B\xi}
         =\ip\eta e ,
\]
where $\ip v e=0$. Hence the targetwise gap has the identity
\begin{equation}\label{eq:second}
 t^\nb_a-\big((Z_v)_{aa}/\pi_a-1\big)
   =\ip\eta e=\ip\xi{L\xi}+\ip\eta{C\eta}\ \ge0 .
\end{equation}
Equality implies $L\xi=C\eta=0$, so $\xi$ is constant and
$e=B\xi+C\eta=0$. Conversely $e=0$ implies equality. Equation
\eqref{eq:variance} and averaging with weights $\pi_a$ give the
claimed criterion.

For the bounds, put $t=QJ(\xi+\eta)$. The equation $L\xi=B^*\eta$
gives $PJ(\xi+\eta)=\xi$, so $J(\xi+\eta)=\xi+t$. Since $J$ is an
isometry and the two summands on either side are orthogonal,
$\|t\|=\|\eta\|$. Write $s=(\eta+t)/2$, $c=(\eta-t)/2$, and let $D$
multiply by $d_i$ on each fibre. Then $s\perp c$ and
\[
 e=(D-2I)\eta+2s,\qquad
 g:=\ip\eta e=\ip\eta{(D-2I)\eta}+2\|s\|^2 .
\]
For each oriented edge with current vertex $i$, weighted Cauchy--Schwarz
gives $((d_i-2)\eta+2s)^2\le d_i((d_i-2)\eta^2+2s^2)$.
Summing yields $\|e\|^2\le\Delta g$. If $\delta>2$, we also have
$g\ge(\delta-2)\|\eta\|^2$. Together with
$g^2\le\|\eta\|^2\|e\|^2$ this gives $(\delta-2)g\le\|e\|^2$.

To handle degree two, observe that $2c=(I-J)(\xi+\eta)$, hence
$Jc=-c$. Also $Pc=0$. The magnitude of this antisymmetric flow is
therefore constant along each degree-two chain. Write $f_H$ for the
restriction to states whose current vertex has degree at least three,
extended by zero. A chain of length $k$ contributes $2k$ oriented edges
and two endpoint states, even when its two endpoints coincide. Thus
$\|c\|^2\le\ell\|c_H\|^2$. Using $c=\eta-s$ gives
\[
 \|\eta\|^2=\|c\|^2+\|s\|^2
 \le 2\ell\|\eta_H\|^2+(2\ell+1)\|s\|^2\le 2\ell g .
\]
The last inequality uses $d_i-2\ge1$ on $H$ and $\ell\ge1$.
It follows that $g\le2\ell\|e\|^2$. Averaging over $a$ in
\eqref{eq:variance} and \eqref{eq:second} yields the stated bounds.
\end{proof}

\paragraph{Formal verification.}
The supplement contains Lean~4.19.0 proofs with mathlib v4.19.0.
For Theorem~\ref{thm:start}, the \texttt{Walks} namespace contains
\texttt{fixed\_target\_comparison} and \texttt{stationary\_limits}.
Together they give the comparison, equality
criterion and convergence of the degree-weighted finite-horizon means.
\texttt{Kemeny.double\_average} gives the strict comparison except on
$K_2$. For Proposition~\ref{prop:theta}, \texttt{Theta.means} verifies
the two formulas by first-step equations at the terminals and at every
interior position. \texttt{Theta.square\_family} proves the ratio limit
along $q=p^2$. The separately encoded 23-vertex graph is checked in
\texttt{Witness.lean}.

In \texttt{Kemeny.lean}, the order and equality criterion of
Theorem~\ref{thm:trace} appear in \texttt{comparison\_\allowbreak and\_\allowbreak equality}.
The two bounds involving the maximum degree and the corridor length appear
in \texttt{variance\_\allowbreak bounds}. The sharper estimate for minimum degree
above two appears in \texttt{min\_\allowbreak degree\_\allowbreak bound}. The proof constructs the inverse defining
$\kp$ and checks that compression gives the same diagonal in the
forward-edge convention. See \texttt{projected\_trace} in \texttt{Forward}.
The supplement includes the package configuration and the kernel's
axiom reports for the main statements.

\paragraph{Exact numerical checks.}
The supplementary Python programs construct the finite graphs from their
edge sets. One solves the absorbing first-step systems by rational
elimination. The other sums the survival probabilities from the killed
transition matrix using Newton's identities and checks every first-step
residual exactly. These checks use only the Python standard library.

\section*{Use of generative AI}
OpenAI Codex (GPT-6) was used to derive the arguments, search for and check
the counterexample, implement the exact verification, investigate earlier
work, formalize the results in Lean, and draft and revise the manuscript.

\begin{thebibliography}{9}
\bibitem{ADGL} K. Altisen, S. Devismes, A. Gerbaud and P. Lafourcade,
Analysis of random walks using tabu lists, in \emph{Structural Information
and Communication Complexity}, LNCS 7355, Springer, 2012, 254--266.
\href{https://doi.org/10.1007/978-3-642-31104-8_22}{Publisher}.
\href{https://www-verimag.imag.fr/PUBLIS/uploads/cizib1712.pdf}{Author version}.
\bibitem{TR} K. Altisen, S. Devismes, A. Gerbaud and P. Lafourcade,
\emph{Comparisons of mean hitting times for tabu random walks on finite
graphs}, Verimag Research Report TR-2012-6, 2012.
\href{https://www-verimag.imag.fr/TR/TR-2012-6.pdf}{Technical report}.
\bibitem{BFG} J. Breen, N. Faught, C. Glover, M. Kempton, A. Knudson and
A. Oveson, Kemeny's constant for nonbacktracking random walks,
\emph{Random Structures Algorithms} \textbf{63} (2023), 343--363.
\href{https://doi.org/10.1002/rsa.21144}{doi:10.1002/rsa.21144}.
\bibitem{BKKS} J. Breen, M. Kempton, A. Knudson and M. Shumway,
On defining Kemeny's constant for non-backtracking random walks,
\emph{Linear Algebra Appl.} (2026).
\href{https://doi.org/10.1016/j.laa.2026.02.034}{doi:10.1016/j.laa.2026.02.034}.
Section and question numbers here follow
\href{https://arxiv.org/html/2510.06650v1}{arXiv:2510.06650v1}.
\bibitem{FTT} D. Fasino, A. Tonetto and F. Tudisco, Hitting times for
second-order random walks, \emph{Eur. J. Appl. Math.} \textbf{34}
(2023), 642--666.
\href{https://doi.org/10.1017/S0956792522000213}{doi:10.1017/S0956792522000213}.
\bibitem{Neal} R. M. Neal, Improving asymptotic variance of MCMC estimators:
non-reversible chains are better, Technical Report 0406, Department of
Statistics, University of Toronto, 2004.
\href{https://arxiv.org/abs/math/0407281}{arXiv:math/0407281}.
\bibitem{OW} R. Ortner and W. Woess, Non-backtracking random walks and
cogrowth of graphs, \emph{Canad. J. Math.} \textbf{59} (2007), 828--844.
\href{https://doi.org/10.4153/CJM-2007-035-1}{doi:10.4153/CJM-2007-035-1}.
\end{thebibliography}
\end{document}
